\subsection{Rules, Principles, and the Role of Mindset}
\label{sec:5.1.1}
In humans, moral development moves from rule-based reasoning to reasoning grounded in general principles. That progression is the model this section adapts for an AI system: a path from a rudimentary grasp of established rules of conduct to sophisticated ethical judgment brought to bear on decisions no rule fully specifies.

\runin{Three kinds of reasoning, and not three ages}
The psychologists Jean Piaget and Lawrence Kohlberg mapped this territory in humans, and their stage theories anchor what follows \autocite{piaget1932moral,kohlberg1969stage}. Reasoning starts at a pre-conventional level, where consequences do the driving — rules get followed to avoid punishment or gain reward. It moves through a conventional stage, where societal norms and expectations get internalized as one's own. It can reach a post-conventional level, where morality answers to principles like justice, fairness, and individual rights, which a given society may or may not happen to expect.

What a system built this way needs from that account is the distinction between the three kinds of reasoning, and nothing about the order in which they arrive. The box below sets out why the sequence will not carry the weight; the distinction will.

A self-driving car makes the three kinds concrete. One system follows traffic law as written and has no resources beyond it: at a blocked lane with a solid line, it stops indefinitely, because crossing the line is prohibited. A second reads the local convention that everyone crosses the line here when the lane is blocked, and does what the road expects. A third can say why the line is painted there, which is what lets it recognize the case where the convention is wrong — the blind curve where local practice is a habit people have survived rather than a rule that protects them — and depart from both the law as written and what the road expects.

These are three capacities and not three ages. A system can have the third without the second, which is the ordinary failure of a machine that reasons impeccably about principle and drives in a way that terrifies everyone around it. Nothing in the design requires passing through them in order, and there is no maturity at the end of the list.

\begin{esbox}
\boxtitle{What the stage theories will and will not carry}
I have just leaned a chapter on Kohlberg, and elsewhere in this book I have been careful to say when a piece of neuroscience or psychology is contested. The same standard applies here, and applying it changes what this chapter is entitled to claim.

Kohlberg's stage model has been under sustained attack for forty years, on three fronts. The first is cross-cultural. A review of forty-five studies across twenty-seven countries found the basic stage sequence broadly supported, but the highest stages largely absent outside urban Western samples — in Kohlberg's own early work in village settings, stages five and six did not appear at all \autocite{snarey1985crosscultural}. Either most of humanity stops short of moral maturity, or the instrument mistakes a particular culture's idiom of justification for the summit of moral development. The second front is the gap between reasoning and doing. A critical review of the literature relating moral judgment to actual conduct — honesty, altruism, delinquency — found the relationship real but far weaker and more conditional than a stage model implies, and the judgment-action gap has been a standing problem for cognitive-developmental accounts ever since \autocite{blasi1980bridging}. The third is Carol Gilligan's, and section~\ref{sec:5.7.1} cites her for care ethics without saying where her argument came from: she was working alongside Kohlberg when she noticed that the women she interviewed reasoned about moral problems in terms of relationship and responsibility, and that his scoring method recorded the difference as a deficit \autocite{gilligan1982different}. Care ethics enters this book as a critique of the framework this section is built on.

So why use it here at all, when section~\ref{sec:2.2.1} declines to build on mirror neurons for less?

Because the two cases fail differently. The mirror-neuron claim was that a specific mechanism — motor simulation — implements a specific capacity, and the argument for building on it depended entirely on the mechanism being what it was said to be. When the mechanism did not hold up, nothing was left to transfer. Kohlberg's stages are not a mechanism. They are a description of an order: rules first, then norms, then principles that can override both. What the criticisms above damage is the claim that this order is universal, that its top rung is where moral development is heading, and that where a person sits on it predicts what they will do. What survives is narrower — that the three ways of reasoning are distinguishable, and that a system which has only the first is missing something the third supplies. That is what this chapter uses.

Two things follow, and I would rather state them here than let the chapter imply otherwise. The order is not a developmental law, so nothing about an AI system's design should assume it must pass through the stages in sequence or that arriving at the third is maturity. And the disanalogy is not bridged by anything I say here: a child moves through these stages under pressures an AI system does not face — peers whose approval it needs, a body, a status that can be lost. Section~\ref{sec:5.2.1} works through what that costs in the case of mortality and guilt, and section~\ref{sec:5.7.1} says the same thing about play. This chapter borrows a vocabulary for distinguishing kinds of moral reasoning. It does not borrow a theory of how a mind comes to have them.

\end{esbox}
\runin{A system that treats its own errors as data}
Carol Dweck's distinction between a fixed mindset, which treats ability as static, and a growth mindset, which treats it as built through effort and experience, matters here for a narrower reason than it usually gets cited for \autocite{dweck2006mindset}. A system that treats its own moral errors as data has somewhere to go after getting one wrong. A system built to defend its prior outputs does not.

That distinction is where chapter~\ref{sec:3}'s hardest question comes to rest, and it is worth stating plainly rather than leaving between the lines. Chapter~\ref{sec:3} asks for a floor that holds, and the word invites a standard nothing has ever met: never crossed, in any circumstance, for the life of the system. No person clears that bar. Nobody who has raised a child expects them to, and a design goal that no example of the thing being modeled has ever satisfied is a goal that has stopped describing anything.

What the argument actually needs is weaker and available. A floor holds, in the sense chapter~\ref{sec:3} can deliver, when the bearer can be shown to have violated it in terms it accepts, so that the violation registers as one and changes what happens next. That is the difference between a system that did something wrong and a system for which nothing counts as doing something wrong. Section~\ref{sec:3.4} asks whether a bearer's account of itself can be believed, and this is the disposition that question is really about: not accuracy of self-report, but whether the system has any stake in the answer being unflattering.

The engineering version is unglamorous. Incremental learning lets a moral framework absorb cases as they arrive rather than freezing at training time. Exposure to unfamiliar situations widens the base a judgment is drawn from. Human correction supplies what section~\ref{sec:5.1.3} calls apprenticeship. Self-examination lets a system find a pattern in its own outputs that no single output reveals. None of those is remarkable on its own, and what they are for is: keeping open the possibility that the system is the one that is wrong.

The second ingredient is the harder one. A system with room to update its judgment still needs something worth updating toward, and the range of views it actually meets is set by whoever assembles them.
